% ch11_c §11.6–11.9 (书页 394–404 / p408–p418)
%--B-TAIL--
%JOIN% 度的变化，即可写出比热的公式。恰在转变温度处，比热有一个数值为
\begin{equation*}
C_{es} = 1.43\,\gamma T_c, \tag{11.50}
\end{equation*}
的跃变，其中 $\gamma$ 是正常态电子比热公式 (4.46) 中的系数。在更低的温度下，能隙趋于在比热中占据主导，但像
\begin{equation*}
C_{es} \approx e^{-\Delta(T)/T} \tag{11.51}
\end{equation*}
这样简单的公式并不适用，除非 $\Delta$ 已经稳定到其极限值 $\Delta_0$。

%FIG{206}: 超导体的比热，与正常金属电子比热的比较。

\section{持续电流}

超导态最引人注目的特征是，金属对于电流的流动似乎完全不表现出任何电阻。要在 BCS 理论中理解这一点，我们注意到能隙并不系结于晶格——它不像半导体中那样固定在倒格子空间的区界上——而只是出现在费米能级处。我们上面所用的论证并不依赖于费米面是球形的；对于任意形状的费米面，它也应当成立。

如 §7.2 所示，电场的作用是使费米分布相对于倒格子整体位移一个恒定矢量。电流 $\mathbf{J}$ 对应于自由电子球的如下位移：
\begin{equation*}
\boldsymbol{\delta}\mathbf{k} = \frac{m}{ne\hbar}\,\mathbf{J} \tag{11.52}
\end{equation*}
能隙并不妨碍这一移动；它随费米面一起被带到新的位置。

要用 BCS 理论的语言描述载流状态，我们只需假设每个电子对的质心移动了 $\boldsymbol{\delta}\mathbf{k}$。在构造式 (11.28) 的准粒子算符 $\beta_{\mathbf{k}}$、$\beta_{\mathbf{k}}^{*}$ 时，我们如今把态 $(\mathbf{k}+\boldsymbol{\delta}\mathbf{k},\uparrow)$ 与态 $(-\mathbf{k}+\boldsymbol{\delta}\mathbf{k},\downarrow)$ 相配；论证的其余部分照旧不变。

%FIG{207}: （a）半导体中，能隙系结于布里渊区。（b）超导体中，能隙由费米面携带。

唯一需要指出的一点是，这样一种电子分布的总动能将增加
\begin{equation*}
\delta\mathcal{E} \approx \frac{\hbar^{2}(\delta k)^{2}}{2m} \tag{11.53}
\end{equation*}
的量（每个电子）。如果这个能量超过了超导态的结合能 (11.38)，我们就不能指望系统仍凝聚在能隙之下。这就为所能承载的电流设定了一个极限。

但载流态最重要的性质是：一旦建立起来，它就不容易被破坏。如 §7.5 所讨论的，阻碍正常电流的那些过程是“单电子”跃迁——单个粒子通过与声子或杂质的相互作用，从费米面的一侧被散射到另一侧。在超导态中，这样的过程只能通过产生准粒子来进行——在电子原来所在处（比如说 $\mathbf{k}+\mathbf{q}$）产生一个空穴，而电子又在 $(-\mathbf{k}+\mathbf{q})$ 处重新出现。这需要能量 $\Delta$，而单个声子提供不了这份能量。要把电流降到零，就得设法让所有电子同时慢下来，使整个位移了的费米面平滑地退回原点。对这样一种可能性极小的过程，并不存在简单的机制，因此在一段环形的超导材料中，环行电流一旦建立，几乎可以永不衰减地持续下去。

BCS 理论似乎关键性地依赖于存在明晰定义的 Bloch 态 $|\mathbf{k},\uparrow\rangle$ 和 $|-\mathbf{k},\downarrow\rangle$。然而，超导电性却常常在极无序、“肮脏”的合金中被观察到；在这些合金里，电子作普通导电时的平均自由程短得使我们无法把 $\mathbf{k}$ 当作好的量子数来处理。真实的情况是：超导“配对”可以在任何两个经由时间反演相互关联的态之间发生（§3.11）。全部必要条件不过是：电子在反转其自旋之后，能够沿原路折返它穿过材料的路径，无论这条路径多么迂回曲折。这样一来，Fröhlich 相互作用的理论看上去会非常复杂，但只要系统的本征态能够按时间反演配对来分类，且费米能处于态密度高的区域，那么在原则上没有任何东西妨碍一个宏观超导凝聚相的出现。基于同样的理由，把近独立电子态重新诠释为多体准粒子激发（§5.8），也不构成 BCS 理论的障碍。

然而，如果材料中含有内建的磁场，例如在铁磁体中，时间反演对称性的假设便不再成立（§9.3），因为我们无法假设在让单个电子沿其路径折返的假想过程中，这些磁场也会一并被反转。因此，相当低的磁性杂质浓度（§10.6）就足以摧毁超导电性。在杂质浓度稍低一些时，则可以观察到一种称为无能隙超导电性（gapless superconductivity）的状态——表面上是超导的、宏观上是有序的（§11.9），但准粒子谱中却没有实际的能隙 (11.40)。

\section{London 方程}

超导体在外磁场中的行为颇为奇特。其中一部分行为表面上可以理解为电导率无限的后果。我们无法在超导体内部建立起磁场，因为在磁场建立时，样品表面会感应出持续电流（persistent current），从而把内部屏蔽起来。但这个解释并不足以说明全部现象。

我们需要考察由具有矢势 $\mathbf{A}(\mathbf{r})$ 的磁场感应出的电流。按照式 (10.3)，这会引起动能算符的改变；在态 $\psi(\mathbf{r})$ 中，这个改变的期望值为
\begin{equation*}
\psi^{*}\left\{\frac{1}{2m}\left(\frac{\hbar}{i}\nabla - \frac{e}{c}\mathbf{A}\right)^{2} - \left(-\frac{\hbar^{2}}{2m}\nabla^{2}\right)\right\}\psi = -\frac{1}{c}\,\mathbf{J}(\mathbf{r})\cdot\mathbf{A}(\mathbf{r}), \tag{11.54}
\end{equation*}
在这里，我们是把这个局部的“动能密度”等同于密度为 $\mathbf{J}(\mathbf{r})$ 的电流在矢势 $\mathbf{A}(\mathbf{r})$ 中运动时相应的经典能量密度。于是，电流密度就是 (11.54) 对 $\mathbf{A}(\mathbf{r})$ 的导数，其形式为
\begin{equation*}
\mathbf{J}(\mathbf{r}) = \frac{e\hbar}{2m}\,\frac{1}{i}\left(\psi^{*}\nabla\psi - \psi\nabla\psi^{*}\right) - \frac{e^{2}}{mc}\,\psi^{*}\psi\,\mathbf{A}. \tag{11.55}
\end{equation*}
第一项就是大家熟知的无磁场时的电流公式。

在普通金属中，静场下这两部分贡献几乎相消，只剩下微弱的抗磁性（参见 §10.1）。然而在超导体中，“波函数” $\psi$ 的那些态（将在 §11.9 中更详细地讨论）具有某种“刚性”。举例来说，假设磁场的施加对波函数毫无影响，$\psi$ 保持等于零场中的值 $\psi_0$，那么我们可以肯定
\begin{align*}
\mathbf{J}_0(\mathbf{r}) &= \frac{e\hbar}{2mi}\left(\psi_0^{*}\nabla\psi_0 - \psi_0\nabla\psi_0^{*}\right)\\
&= 0. \tag{11.56}
\end{align*}
于是在超导体中，剩下的便是
\begin{align*}
\mathbf{J}(\mathbf{r}) &= -\frac{e^{2}}{mc}\,\psi^{*}(\mathbf{r})\psi(\mathbf{r})\,\mathbf{A}(\mathbf{r})\\
&= -\frac{e^{2}n}{mc}\,\mathbf{A}(\mathbf{r}), \tag{11.57}
\end{align*}
其中 $n$ 是当地的电子密度。

这个不再出现假想“波函数”的公式，就是著名的 London 方程（London equation）；它对解释超导体的许多性质大有帮助。例如，我们知道电导率是无限的，因此样品内部没有电场。静磁场将通过 Maxwell 方程
\begin{equation*}
\nabla\wedge\mathbf{H} = \frac{4\pi}{c}\,\mathbf{J}. \tag{11.58}
\end{equation*}
与电流 $\mathbf{J}$ 相联系。但 (11.57) 可以写成
\begin{align*}
\mathbf{H} &= \nabla\wedge\mathbf{A}\\
&= -\frac{mc}{e^{2}n}\,\nabla\wedge\mathbf{J}. \tag{11.59}
\end{align*}
把 (11.58) 与 (11.59) 结合起来，我们得到
\begin{equation*}
\nabla^{2}\mathbf{H} = \frac{1}{\lambda^{2}}\,\mathbf{H}, \qquad \nabla^{2}\mathbf{J} = \frac{1}{\lambda^{2}}\,\mathbf{J}, \tag{11.60}
\end{equation*}
其中
\begin{equation*}
\lambda = \sqrt{\frac{mc^{2}}{4\pi ne^{2}}}. \tag{11.61}
\end{equation*}
对于恰在超导体表面之外并平行于表面的磁场 $\mathbf{H}_0$，(11.60) 的解是
\begin{equation*}
\mathbf{H} = \mathbf{H}_0\,e^{-z/\lambda}, \tag{11.62}
\end{equation*}
其中 $z$ 是深入样品内部量度的距离。这样，磁场迅速衰减，只能穿透 $\lambda$ 的距离——这就是所谓的穿透深度（penetration depth），其量级为 $5\times10^{-6}$ cm。

这一结果并不依赖于磁场是如何建立的。对于在磁场中从转变温度之上冷却下来的样品，它同样成立。当金属变为超导体时，磁场被排了出去。这就是 Meissner 效应（Meissner effect），它无法仅凭无限电导率来解释。

超导态要求时间反演对称性，而且是“完全抗磁的”，这一事实提示：足够强的磁场应当能够改变这个态。确实，当施加的磁场超过临界值 $H_c$ 时，块体金属中的超导电性即被破坏。对第一类超导体（type I superconductor，见 §11.10），这个场可以由比热曲线通过热力学论证计算出来。临界场处的场能 $H_c^{2}/8\pi$ 必须等于超导态与正常态之间的自由能之差。举例来说，假设超导态的比热正比于 $T^{3}$，而正常态具有普通的电子比热 $\gamma T$，那么可以证明临界场应是温度的函数：
\begin{equation*}
H_c = \sqrt{(2\pi)\gamma T_c\{1-(T/T_c)^{2}\}}, \tag{11.63}
\end{equation*}
其中 $T_c$ 是临界温度 (11.48)。实际上，正如在 (11.51) 处所指出的，超导体的比热是温度的复杂函数，所以这个公式只是近似正确。

值得注意的是，这一论证假设我们能够把\emph{全部}磁场排除在样品之外。对于厚度与 $\lambda$ 相当的很小或很薄的金属块，磁场会穿透进去，热力学论证便不再有效。这时超导态可以存在于比 (11.63) 算出的值高得多的磁场之下。第二类超导体（type II superconductor）中发生的正是这种情形：在那里，即使在块体样品中，这类细丝也可以通过与位错、杂质等的相互作用而稳定下来（见 §11.10）。

\section{相干长度}

推导 London 方程 (11.57) 时的假设是：超导波函数是“刚性的”，不受磁场施加的影响。这显然不正确；我们知道，足够强的磁场会破坏超导态。我们需要计算 (11.55) 中的波函数 $\psi$ 在微扰 $\mathbf{A}(\mathbf{r})$ 作用于超导基态时所产生的改变。

这在 §11.3 的 Bogoliubov 形式框架内并不难算出。假设我们考虑矢势的一个单独的 Fourier 分量
\begin{equation*}
\mathbf{A}(\mathbf{r}) = \mathbf{A}_{\mathbf{q}}\,e^{i\mathbf{q}\cdot\mathbf{r}}. \tag{11.64}
\end{equation*}
在 $\exp(i\mathbf{k}\cdot\mathbf{r})$ 和 $\exp(i\mathbf{k}'\cdot\mathbf{r})$ 这两个普通自由电子态之间，这一项之一的矩阵元将是
\begin{equation*}
\int e^{-i\mathbf{k}'\cdot\mathbf{r}}\,\mathbf{A}(\mathbf{r})\,e^{i\mathbf{k}\cdot\mathbf{r}}\,\mathrm{d}\mathbf{r} = \mathbf{A}_{\mathbf{q}}\,\delta(\mathbf{k}+\mathbf{q}-\mathbf{k}'). \tag{11.65}
\end{equation*}
这样，用湮没和产生算符 (11.24) 的语言来说，$\mathbf{A}(\mathbf{r})$ 的行为就如同算符
\begin{equation*}
\mathbf{A} = \mathbf{A}_{\mathbf{q}}\,b_{\mathbf{k}+\mathbf{q}}^{*}b_{\mathbf{k}} \tag{11.66}
\end{equation*}
一样。精确到 $\mathbf{A}$ 的一阶，磁场在 Hamiltonian 中产生的实际微扰就是算符
\begin{equation*}
-\frac{e\hbar}{2mci}\left(\nabla\cdot\mathbf{A}+\mathbf{A}\cdot\nabla\right) = -\frac{e\hbar}{2mc}\sum_{\mathbf{k}}\left\{(\mathbf{k}+\mathbf{q})+\mathbf{k}\right\}\cdot\mathbf{A}_{\mathbf{q}}\,b_{\mathbf{k}+\mathbf{q}}^{*}b_{\mathbf{k}}. \tag{11.67}
\end{equation*}
这很容易通过与 (11.65) 和 (11.66) 的类比得到；在所有“裸电子”态之间，右边的算符与左边的算符具有相同的矩阵元。

现在对这个表达式施加变换 (11.28)。用准粒子算符 $\beta_{\mathbf{k}}$、$\beta_{\mathbf{k}}^{*}$ 表出的结果有点复杂，但由于我们是要把微扰作用到超导基态 $|0\rangle$ 上，可以利用 (11.32) 消去所有含有湮没算符的项。剩下的就是含有两个产生算符之积的项；我们的微扰变为
\begin{equation*}
-\frac{e\hbar}{2mc}\sum_{\mathbf{k}}\,(2\mathbf{k}+\mathbf{q})\cdot\mathbf{A}_{\mathbf{q}}\,(v_{\mathbf{k}+\mathbf{q}}u_{\mathbf{k}} - v_{\mathbf{k}}u_{\mathbf{k}+\mathbf{q}})\,\beta_{\mathbf{k}+\mathbf{q}}^{*}\beta_{-\mathbf{k}}^{*}, \tag{11.68}
\end{equation*}
其中 $v_{\mathbf{k}}$、$u_{\mathbf{k}}$ 等是 (11.28) 中的系数。磁场产生净动量为 $\mathbf{q}$ 的准粒子对，正如我们从关于持续电流性质的论证（§11.6）中所应预期的那样。

我们用 (11.68) 作为微扰来计算基态波函数的改变，然后用这个“非刚性”的态代替 $\psi$，重新由 (11.55) 计算电流。结果如下：
\begin{equation*}
\mathbf{J}_{\mathbf{q}} = -\frac{e^{2}n}{mc}\,\mathbf{A}_{\mathbf{q}} + \frac{e^{2}}{2m^{2}c}\sum_{\mathbf{k}}\frac{(2\mathbf{k}+\mathbf{q})\left\{(2\mathbf{k}+\mathbf{q})\cdot\mathbf{A}_{\mathbf{q}}\right\}(v_{\mathbf{k}}u_{\mathbf{k}+\mathbf{q}} - v_{\mathbf{k}+\mathbf{q}}u_{\mathbf{k}})}{\epsilon(\mathbf{k})+\epsilon(\mathbf{k}+\mathbf{q})}, \tag{11.69}
\end{equation*}
其中准粒子对的能量 $\epsilon(\mathbf{k})+\epsilon(\mathbf{k}+\mathbf{q})$ 由 (11.40) 算出。

这一计算中存在一些形式上的困难——关于矢势规范选取的困难，这类困难在磁场中粒子的量子理论中经常出现。这里我们可以假定 $\mathbf{A}(\mathbf{r})$ 是纯横向的，即 $\mathbf{q}\cdot\mathbf{A}_{\mathbf{q}}=0$，从而避开这些困难。为方便起见，把 $v_{\mathbf{k}}$ 等的乘积对称化；我们得到
\begin{equation*}
\mathbf{J}_{\mathbf{q}} = -\frac{e^{2}n}{mc}\,\mathbf{A}_{\mathbf{q}} + \frac{2e^{2}}{m^{2}c}\sum_{\mathbf{k}}\frac{\mathbf{k}(\mathbf{k}\cdot\mathbf{A}_{\mathbf{q}})\left(v_{\mathbf{k}-\frac{1}{2}\mathbf{q}}u_{\mathbf{k}+\frac{1}{2}\mathbf{q}} - v_{\mathbf{k}+\frac{1}{2}\mathbf{q}}u_{\mathbf{k}-\frac{1}{2}\mathbf{q}}\right)}{\epsilon(\mathbf{k}-\frac{1}{2}\mathbf{q})+\epsilon(\mathbf{k}+\frac{1}{2}\mathbf{q})}. \tag{11.70}
\end{equation*}
利用 (11.36) 中 $u_{\mathbf{k}}$ 和 $v_{\mathbf{k}}$ 用能量 $\xi(\mathbf{k})$ 表出的公式，这个和式很容易算出。我们不打算详细追述这一论证；它不过是初等的几何。我们只是把 (11.70) 写成
\begin{equation*}
\mathbf{J}_{\mathbf{q}} = -\frac{c}{4\pi}\,\Gamma(\mathbf{q})\,\mathbf{A}_{\mathbf{q}} \tag{11.71}
\end{equation*}
的形式，用以表示电流与矢势之间的关系；$\Gamma(\mathbf{q})$ 是所施加磁场的波数 $q$ 的一个完全确定的函数。

从 (11.70) 的形式可以相当清楚地看出，当 $q\to 0$ 时和式消失。因此，对缓慢变化的磁场，London 公式 (11.57) 成立；一般说来，Meissner 效应随之而来。但随着 $q$ 增大，(11.70) 中的和式——它本质上是正的——趋于抵消首项。的确，对大的 $q$ 值，我们考虑的是能量远大于能隙的激发，准粒子的行为就像普通的电子和空穴（依其在费米能级之上还是之下而定）。这时可以证明——这是运用该理论的一个很好的练习——London 项恰好被抵消。换言之，当 $q$ 变大时，
\begin{equation*}
\Gamma(q) \to 0 \tag{11.72}
\end{equation*}
$\Gamma(q)$ 在这两个极限之间的实际形式如图 208 所示。很自然可以设想（而且这同样可以由 (11.70) 验证），当 $q$ 大到足以产生一个激发时——也就是当 $q$ 超过
\begin{equation*}
\hbar v_F q_c \sim \Delta_0. \tag{11.73}
\end{equation*}
所给出的 $q_c$ 时——我们将从“London”区制过渡到 (11.72) 所示的区制。

%FIG{208}: 响应函数：（a）在波数空间中，（b）在真实空间中。

这个公式所定义的量 $1/q_c$，就是在纯金属的理想情形下\emph{相干长度}（coherence length）$\xi$ 的值。

为了看出它的意义，让我们把 (11.71) 表成\emph{真实空间}中电流与矢势之间的关系。它变为
\begin{equation*}
\mathbf{J}(\mathbf{r}) = \int \Gamma(\mathbf{r}-\mathbf{r}')\,\mathbf{A}(\mathbf{r}')\,\mathrm{d}\mathbf{r}', \tag{11.74}
\end{equation*}
其中 $\Gamma(\mathbf{r})$ 是 $\Gamma(\mathbf{q})$ 的 Fourier 变换。在 London 理论中 $\Gamma(\mathbf{q})$ 是常数，所以 $\Gamma(\mathbf{r})$ 是 $\delta$ 函数；$\mathbf{J}(\mathbf{r})$ 的值只依赖于 $\mathbf{A}(\mathbf{r})$ 的\emph{局域}值，如 (11.57) 那样。但 (11.71) 的变换要更复杂。$\Gamma(\mathbf{r})$ 的范围将是相干长度 $\xi$，于是我们将得到 $\mathbf{J}$ 与 $\mathbf{A}$ 之间的一个\emph{非局域}关系。它与反常趋肤效应理论中使用的 Chambers 公式 (8.101) 颇为相似。

这一由 Pippard 从唯象上建立的结果，对于反常趋肤效应、磁场穿透等性质的细致分析十分重要。实际上，$\xi$ 就是一个 Cooper 对的“尺度”；它是配对的电子的相位必须发生关联的距离。在“肮脏”的金属中，当电子平均自由程 $\Lambda$ 小于 $1/q_c$ 时，相干长度 $\xi$ 的表观值不会超过 $\Lambda$ 本身。这时相位关联在一定程度上被破坏，但配对本身并未被摧毁，材料仍是超导体。

有意思的是看一看这种相干性在数学理论中是如何表达的。从 (11.67) 到 (11.68) 的步骤把它显示得非常清楚。算符 (11.67) 引起单电子态之间的散射。在正常金属中，我们应当把这些事件当作相互独立的事件来处理，并把各个矩阵元的平方相加。这样，在 (11.69) 中我们将恰好有对应于
\begin{equation*}
(v_{\mathbf{k}}u_{\mathbf{k}+\mathbf{q}})^{2} + (u_{\mathbf{k}}v_{\mathbf{k}+\mathbf{q}})^{2} \tag{11.75}
\end{equation*}
的项出现在分子中。

但在超导态中，该算符以\emph{单一}的矩阵元
\begin{equation*}
(v_{\mathbf{k}}u_{\mathbf{k}+\mathbf{q}} - u_{\mathbf{k}}v_{\mathbf{k}+\mathbf{q}}) \tag{11.76}
\end{equation*}
产生两个动量近乎相反的激发，再取其平方。这给出的结果与 (11.75) 不同；正是乘积项在微扰所产生的两个准粒子之间引起关联，即相干性。但这一项只对小于 $q_c$ 的 $q$ 才重要；这正是相干长度的意义所在。

\section{非对角长程序}

超流电流可以沿着长达许多千米的回路流动。这样的回路上某一点的物理条件会受到相距极远的其他各点条件的影响。因此，电子系统表现出一种特殊形式的\emph{长程序}（long range order），我们想从数学上刻画它，同时适当地顾及当我们从一个区域走到另一区域时材料、温度或磁场的改变。

在半导体的宏观理论中，这些局域特征被概括为各类载流子的密度和迁移率。类似地，在 Gorter--Casimir 二流体模型（two-fluid model）中，我们假设超流电流由数密度为 $n_s$ 的“超流电子”携带，$n_s$ 依赖于温度等。如果每个这样的电子以速度 $\mathbf{v}_s$ 运动、携带电荷 $e^{*}$（不一定是普通的电子电荷 $e$），那么总的超流电流密度必须是
\begin{equation*}
\mathbf{J}_s = n_s e^{*}\mathbf{v}_s. \tag{11.77}
\end{equation*}
这个模型在 BCS 理论的框架内几乎可以得到论证：把“正常”电子的平衡浓度
\begin{equation*}
n_n = n - n_s, \tag{11.78}
\end{equation*}
认同为 (11.41) 所给出的准粒子激发的浓度。在 $T=0$ 时没有准粒子，故 $n_s=n$；在临界温度 $T_c$ 处，$n_s$ 迅速降到零。我们甚至可以在热力学上逼近图 206 的比热反常，办法是在超流相与“正常”相之间指定一个自由能差。于是我们可以设想，例如，超导隧穿（§11.11）这类空间性的效应，在数学上可以用 $n_s$ 的局域值随位置的变化来描述。

然而这个描述并不充分。正如我们在 §11.8 中看到的，超导基态的特征是：粒子对的波函数在大于 (11.73) 所定义的相干长度 $\xi$ 的距离上存在强关联。因此，让我们回到 §11.7 的 London 理论，并且有意为电子密度的超流分量引入一个\emph{宏观波函数} $\Psi(\mathbf{r})$（通常称为 Ginzburg--Landau 序参量（order parameter））。正如我们在 §11.8 中学到的那样，只要我们不涉及在短于相干长度的距离上迅速变化的性质，这可以很好地充当超流序的一个适当的局域量度。

按照量子力学的正则形式，我们自然假设局域超流密度由
\begin{equation*}
n_s(\mathbf{r}) = |\Psi(\mathbf{r})|^{2} \tag{11.79}
\end{equation*}
给出，并且
\begin{equation*}
\mathbf{J}_s(\mathbf{r}) = \frac{e^{*}}{2m^{*}}\left[\Psi^{*}\left\{\frac{\hbar}{i}\nabla - \frac{e^{*}}{c}\mathbf{A}\right\}\Psi + \Psi\left\{-\frac{\hbar}{i}\nabla - \frac{e^{*}}{c}\mathbf{A}\right\}\Psi^{*}\right] \tag{11.80}
\end{equation*}
定义超流电流密度，正如 (10.3) 和 (11.55) 那样。但由于怀疑超流电流的载体是“对”而不是单个电子，我们让 $e^{*}$ 和 $m^{*}$ 保持未定义。

如果 $\Psi(\mathbf{r})$ 处处为实，那么它就可以被消去而换成 $n_s$：(11.79) 和 (11.80) 立即约化为 London 方程 (11.57)，只是把 $e$、$m$ 和 $n$ “重整化”为 $e^{*}$、$m^{*}$ 和 $n_s$。在 §11.8 所讨论的限制之内，这就为 Meissner 效应提供了一个良好的定性描述，其中包括隐含在 $n_s$ 随 $T$ 变化之中的穿透深度对温度的依赖。

但是，除非 $\Psi(\mathbf{r})$ 是\emph{复}的，我们就不能描述均匀材料中的超流电流。这一唯象理论出乎意料的特征是：物理情况直接依赖于序参量的\emph{相位}的变化，而不只是依赖于它的绝对大小：用矩阵力学的语言说，我们必须与非对角长程序（off-diagonal long range order）打交道。

例如，假设我们有一种温度恒定且均匀的材料，其中 $n_s(\mathbf{r})$ 为常数。写成
\begin{equation*}
\Psi(\mathbf{r}) = (n_s)^{\frac{1}{2}}\,e^{i\chi(\mathbf{r})}, \tag{11.81}
\end{equation*}
我们便得到
\begin{equation*}
m^{*}\mathbf{v}_s(\mathbf{r}) = \hbar\nabla\chi(\mathbf{r}) - (e^{*}/c)\,\mathbf{A}(\mathbf{r}) \tag{11.82}
\end{equation*}
在不存在磁场时，相位函数 $\chi(\mathbf{r})$ 扮演 (11.77) 中已定义的超流速度 $\mathbf{v}_s$ 的速度势的角色。但如果 $\chi$ 不随空间变化，就不会有超流电流。

%FIG{209}: 用于磁通量子化（flux quantization）的积分路径。

现在考虑一个被磁场（具有矢势 $\mathbf{A}(\mathbf{r})$）穿过的超导材料回路（图 209）。让我们沿环绕此电路的一条路径积分 (11.82)，路径保持在超导介质内部足够深处，以避开可能发生抗磁屏蔽电流的穿透区。在电路中不存在其他电动势源的情况下，$\mathbf{v}_s(\mathbf{r})$ 沿此

